% =====================================================================
%  ArduPlane Controller / Gain Tuning & Gain Scheduling
%  for a custom uSTOL FDM in SITL (Software-in-the-Loop)
%
%  HOW TO COMPILE (any of these):
%     pdflatex ArduPlane_Tuning_Guide.tex   (run twice for the table of contents)
%   or
%     latexmk -pdf ArduPlane_Tuning_Guide.tex
%
%  Requires a standard TeX Live / MiKTeX install. No special packages.
% =====================================================================
\documentclass[11pt,a4paper]{article}

\usepackage[T1]{fontenc}
\usepackage[utf8]{inputenc}
\usepackage{lmodern}
\usepackage[margin=2cm]{geometry}
\usepackage{longtable}
\usepackage{booktabs}
\usepackage{array}
\usepackage{xcolor}
\usepackage{enumitem}
\usepackage{fancyhdr}
\usepackage{parskip}
\usepackage{listings}
\usepackage{underscore} % lets RLL_RATE_P be typed with literal underscores
\usepackage{url}

\definecolor{codebg}{RGB}{245,245,245}
\definecolor{rulecol}{RGB}{120,120,120}

\lstset{
  basicstyle=\ttfamily\small,
  backgroundcolor=\color{codebg},
  breaklines=true,
  frame=single,
  rulecolor=\color{rulecol},
  columns=fullflexible,
  keepspaces=true,
  showstringspaces=false,
  xleftmargin=4pt,
  framexleftmargin=4pt
}

\newcommand{\pp}[1]{\texttt{#1}} % parameter name

\pagestyle{fancy}
\fancyhf{}
\lhead{\small ArduPlane Tuning Guide --- custom uSTOL FDM}
\rhead{\small\thepage}
\renewcommand{\headrulewidth}{0.4pt}

\title{\textbf{ArduPlane Controller, Gain Tuning \& Gain Scheduling}\\[4pt]
\large for a custom uSTOL FDM in SITL}
\author{Prepared as a teaching reference}
\date{June 2026}

\begin{document}

\maketitle

\noindent This guide covers controller tuning, gain tuning, and gain scheduling
for fixed-wing \textbf{ArduPlane}, written for a \emph{custom flight-dynamics
model} (the \pp{libraries/SITL/LAT\_SIM\_*} files) modelling a uSTOL /
distributed-propulsion aircraft running in SITL. It is built entirely from the
official ArduPilot Plane documentation. The aircraft here is pure fixed-wing
(\pp{Q\_ENABLE,0}), with nine throttle channels (distributed propulsion) plus
conventional aileron / elevator / rudder / flaps. Part~9 shows how to do
everything \textbf{with Mission Planner}; Part~10 shows the same \textbf{without
Mission Planner} (headless / MAVProxy), which suits a SITL workflow.

\tableofcontents
\newpage

% =====================================================================
\section{Part 0 --- The mental model}
% =====================================================================

ArduPlane control is a \textbf{cascade of nested loops}. Outer loops are slow and
set \emph{goals}; inner loops are fast and \emph{achieve} them. You always tune
\textbf{inner-to-outer}.

\begin{lstlisting}
 GPS waypoint  --> L1 Nav ------> roll-angle demand --+
                                                      +--> ATTITUDE --> RATE demand --> RATE PID-FF --> servo
 Altitude/speed --> TECS ------> pitch + throttle ----+                                                (ail/elev/rudder)
\end{lstlisting}

\begin{longtable}{@{}p{3.1cm}p{6.0cm}p{4.6cm}p{1.7cm}@{}}
\toprule
\textbf{Layer} & \textbf{What it decides} & \textbf{Main parameters} & \textbf{Order}\\
\midrule
\endhead
Rate (innermost) & Surface deflection to hit a demanded rate (deg/s) & \pp{RLL\_RATE\_*}, \pp{PTCH\_RATE\_*}, \pp{YAW\_RATE\_*} & 1st\\
Attitude & deg/s needed to reach a demanded angle & \pp{*2SRV\_TCONST}, \pp{*2SRV\_RMAX*} & with rate\\
Speed/Height (TECS) & pitch + throttle for altitude \& airspeed & \pp{TECS\_*}, \pp{THR\_*}, \pp{PTCH\_LIM\_*} & 2nd\\
Navigation (L1) & bank angle to follow the ground track & \pp{NAVL1\_PERIOD}, \pp{NAVL1\_DAMPING}, \pp{WP\_RADIUS} & 3rd\\
\bottomrule
\end{longtable}

\textbf{Why this order matters:} TECS asks the pitch loop for a pitch angle. If
the pitch loop is sloppy, TECS looks broken even though the bug is one layer
down. Same with L1 $\rightarrow$ roll. Never tune an outer loop until the inner
one is solid.

\subsection*{The PID-FF building block (identical on roll, pitch, yaw rate loops)}
\begin{itemize}[leftmargin=1.4em]
  \item \textbf{FF (feedforward)} --- the dominant term for planes. Maps the
  \emph{demanded rate} straight to surface deflection. A well-set FF flies the
  plane; PID just cleans up the error. (This is why planes are FF-led and
  multirotors are P-led.)
  \item \textbf{P} --- reacts to current rate \emph{error}. Too high $\rightarrow$
  overshoot / porpoising.
  \item \textbf{I} --- integrates persistent error (trim, CG offset, wind). Acts
  like auto-trim.
  \item \textbf{D} --- damping on rate-of-change. Too high $\rightarrow$ amplifies
  noise, heats servos, buzzes.
  \item Plus limiters/filters: \pp{IMAX}, \pp{FLTT}, \pp{FLTE}, \pp{FLTD},
  \pp{SMAX} (Part~2).
\end{itemize}

% =====================================================================
\section{Part 1 --- Prerequisites (before touching a gain)}
% =====================================================================

Tuning a bad setup just bakes the badness into your gains.

\begin{enumerate}[leftmargin=1.6em]
  \item \textbf{Control-surface signs.} In \textbf{FBWA}, nose-up $\rightarrow$
  elevator down; roll right $\rightarrow$ right aileron up; etc. In a custom FDM
  this doubly matters: a sign error in
  \pp{LAT\_SIM\_Forces\_and\_moments\_ctrl\_srfce.cpp} \emph{looks} like a tuning
  problem but is a model bug. Verify the servo outputs in SITL before blaming
  gains.
  \item \textbf{Set \pp{SCALING\_SPEED} first, and never move it later.} Set it
  to cruise. It is the anchor for gain scheduling (Part~3). Changing it after
  tuning rescales every PID and corrupts the tune.
  \item \textbf{Airspeed.} In SITL airspeed is synthesised from your FDM, so it is
  effectively perfect --- but the envelope (\pp{AIRSPEED\_MIN}, \pp{AIRSPEED\_MAX},
  \pp{AIRSPEED\_CRUISE}) must match what the FDM can actually do. If the model
  stalls at 22 m/s, \pp{AIRSPEED\_MIN}=20 is dangerous.
  \item \textbf{CG / trim.} Bench-trim neutral. In FBWA at cruise the surfaces
  should sit near-neutral ($<10\%$ throw). If not, it is CG or level-calibration.
  \item \textbf{Servo rate limits.} Set \pp{RLL\_RATE\_SMAX} / \pp{PTCH\_RATE\_SMAX}
  to the servo slew capability. Your FDM has \pp{LAT\_SIM\_servo\_dynamics.cpp};
  make SMAX consistent with the servo model's rate or the limiter and the model
  will fight.
\end{enumerate}

% =====================================================================
\section{Part 2 --- Controller \& gain tuning (the core)}
% =====================================================================

\subsection{The rate controller (firmware 4.1+)}

Per axis (\pp{RLL\_}, \pp{PTCH\_}, \pp{YAW\_}). The right-hand column shows this
aircraft's current roll / pitch values.

\begin{longtable}{@{}p{2.7cm}p{4.7cm}p{5.4cm}p{2.4cm}@{}}
\toprule
\textbf{Param} & \textbf{Meaning} & \textbf{How to set} & \textbf{Value (R/P)}\\
\midrule
\endhead
\pp{\_RATE\_FF}   & Feedforward, demand$\rightarrow$surface & \textbf{Tune first}; from logs, scale until commanded surface $\approx$ needed & 0.400 / 0.660\\
\pp{\_RATE\_P}    & Proportional on rate error & Raise 0.01 steps until overshoot, back off 25--50\% & 0.214 / 0.516\\
\pp{\_RATE\_I}    & Integral & Start $\approx$ FF (approx 1 s correction) & 0.214 / 0.660\\
\pp{\_RATE\_D}    & Derivative damping & Raise 0.001 steps until oscillation, then \textbf{halve} & 0.034 / 0.020\\
\pp{\_RATE\_IMAX} & Integrator clamp & Default 0.666 (two-thirds throw) & 0.666\\
\pp{\_RATE\_FLTT} & Target filter (Hz) & approx loop bandwidth; lower for slow airframe & 3.18 / 2.12\\
\pp{\_RATE\_FLTE} & Error filter (Hz) & Usually 0 (off) & 0\\
\pp{\_RATE\_FLTD} & D-term filter (Hz) & Knock down noise feeding D & 10 / 10\\
\pp{\_RATE\_SMAX} & Slew limit; auto-degains if exceeded & Servo rate capability & 150\\
\bottomrule
\end{longtable}

Attitude (angle $\rightarrow$ rate) layer:

\begin{longtable}{@{}p{4.6cm}p{8.7cm}p{2.2cm}@{}}
\toprule
\textbf{Param} & \textbf{Meaning} & \textbf{Value}\\
\midrule
\endhead
\pp{RLL2SRV\_TCONST} / \pp{PTCH2SRV\_TCONST} & Time constant to reach demanded angle (smaller = snappier) & 0.5 / 0.75\\
\pp{RLL2SRV\_RMAX} / \pp{PTCH2SRV\_RMAX\_UP/DN} & Max rate limits. Pitch trick: $560/V$ (m/s) $\approx \pm 1$ g & 75 / 75,75\\
\pp{PTCH2SRV\_RLL} & Pitch-up added in turns to hold altitude. Default 1.0, range 0.8--1.3 & 1.0\\
\bottomrule
\end{longtable}

\textbf{Manual tuning sequence (canonical order):}
\begin{enumerate}[leftmargin=1.6em]
  \item \textbf{FF first.} Fly FBWA, log, plot demanded rate vs achieved surface,
  scale FF so they match. FF doing most of the work is the goal.
  \item \textbf{I = FF} as a starting point.
  \item \textbf{P} up in 0.01 steps until overshoot/porpoising $\rightarrow$ back off.
  \item \textbf{D} up in 0.001 steps until it oscillates $\rightarrow$ halve.
  \item Re-touch P slightly if D let you go higher.
\end{enumerate}

This aircraft's tune already shows the fingerprints of this method:
\pp{RLL\_RATE\_I} = \pp{RLL\_RATE\_P} = 0.214, and \pp{PTCH\_RATE\_I} =
\pp{PTCH\_RATE\_FF} = 0.660 --- exactly the ``I approx FF'' rule and an AUTOTUNE
result.

\subsection{Yaw (damper + sideslip)}
\begin{itemize}[leftmargin=1.4em]
  \item \pp{YAW2SRV\_DAMP} --- raise 0.05 until tail-wag, then halve. Tune first.
  \item \pp{YAW2SRV\_RLL} --- turn coordination, default 1.0 (range 0.7--1.4).
  \item \pp{YAW2SRV\_SLIP} / \pp{YAW2SRV\_INT} --- sideslip control; only useful if
  the FDM has meaningful side-force. Slender/flying-wing $\rightarrow$ skip slip,
  keep damper.
  \item \pp{KFF\_RDDRMIX} --- feedforward rudder during rolls to kill adverse yaw
  ($\leq 1.0$).
  \item Rate-based yaw (4.2+): \pp{YAW\_RATE\_ENABLE}=1 plus \pp{YAW\_RATE\_P/I/D/FF},
  autotunable.
\end{itemize}

\subsection{AUTOTUNE --- the recommended path (works great in SITL)}

AUTOTUNE flies like FBWA but learns roll/pitch (and yaw if enabled) FF/P/I/D plus
slew limits from your stick inputs.

\begin{longtable}{@{}p{1.6cm}p{2.6cm}p{10.5cm}@{}}
\toprule
\textbf{Level} & \textbf{Target rate} & \textbf{Use}\\
\midrule
\endhead
2--3 & 30--40 deg/s & slow gliders\\
6 (default) & 75 deg/s & beginner--intermediate, good default (this aircraft)\\
7 & 90 deg/s & sharper, experienced\\
8 & 120 deg/s & light foam under 2 kg\\
9--11 & 160--300 deg/s & advanced, post-initial only\\
0 & --- & special: keeps your rates/TCONST, only optimizes FF/P/I/D\\
\bottomrule
\end{longtable}

Yaw is always tuned to 90 deg/s. Pick the level by what the aircraft can
\emph{physically} roll/pitch --- overshooting capability gives a bad tune.

\begin{itemize}[leftmargin=1.4em]
  \item \textbf{How to fly it:} put it on a switch (\pp{RCx\_OPTION}=107) or a
  flight mode, then make sharp stick doublets --- tuning only engages when the
  demanded rate exceeds 40\% of the level's max. Do at least 20 inputs per axis,
  finish all axes before disarming. P and D save per-axis as you go; FF and I save
  on mode exit. Bailing mid-axis reverts that axis.
  \item \textbf{Compare:} \pp{RCx\_OPTION}=180 toggles original vs tuned gains in
  any mode to A/B them.
  \item \textbf{Light / no-airspeed edge cases:} \pp{FLIGHT\_OPTIONS} bit 6 (+64)
  stops PID inflation on takeoff without airspeed; raise \pp{SCHED\_LOOP\_RATE} to
  200 Hz for very light/fast craft.
\end{itemize}

\textbf{In SITL specifically:} AUTOTUNE is the fastest way to re-baseline a tune
after \emph{any} FDM revision. Every aero change (e.g.\ ``\pp{Cmq} updated'',
``post-stall \pp{CL} added'') shifts the right gains; a 2-minute SITL AUTOTUNE
re-baselines instantly, with no servo to burn out.

% =====================================================================
\section{Part 3 --- Gain scheduling (airspeed speed-scaling)}
% =====================================================================

\subsection{The fundamental problem}
Control-surface effectiveness scales with dynamic pressure, proportional to
$V^2$. A deflection that gives 30 deg/s at 20 m/s gives roughly $4\times$ the
moment at 40 m/s. A fixed gain that is stable at cruise will be sluggish at min
speed and oscillate at max speed. So gains \emph{must} change with speed.

\subsection{ArduPlane's primary mechanism: automatic airspeed scaling}
ArduPlane does this continuously, on the rate-loop output:

\begin{lstlisting}
 scale     = (SCALING_SPEED / current_airspeed)^2     (clamped)
 servo_out = scale * PID_FF_output
\end{lstlisting}

\begin{itemize}[leftmargin=1.4em]
  \item At \textbf{high speed} $\rightarrow$ scale $<1$ $\rightarrow$ surfaces move
  \textbf{less}.
  \item At \textbf{low speed} $\rightarrow$ scale $>1$ $\rightarrow$ surfaces move
  \textbf{more}.
  \item \pp{SCALING\_SPEED} is the anchor where scale $=1$. \textbf{You tune at
  this speed}, and the controller extrapolates everywhere else.
  \item Works \textbf{with or without} an airspeed sensor: with \pp{ARSPD\_USE}=1
  it uses measured speed; without, it uses the EKF estimated speed (GPS+IMU). In
  SITL the FDM feeds a clean airspeed, so scheduling is essentially ideal.
  \item The scaling range is bounded by \pp{AIRSPEED\_MIN}/\pp{AIRSPEED\_MAX}.
  Rule of thumb: \pp{AIRSPEED\_MAX} $\approx 2\times$ cruise, \pp{AIRSPEED\_MIN}
  $\geq 1.25\times$ stall. The factor is clamped so you do not get insane gains
  near zero speed.
\end{itemize}

\textbf{The critical footgun:} because scheduling keys off \pp{SCALING\_SPEED},
changing it after tuning silently rescales all PIDs and ruins the tune. If you
change cruise speed, revert PIDs to default and retune at the new speed.

\subsection{Other adaptive (scheduling-like) mechanisms}
\begin{itemize}[leftmargin=1.4em]
  \item \pp{*\_RATE\_SMAX} (slew-rate limiting) --- if surface slew exceeds SMAX,
  ArduPilot automatically reduces effective gain until back in limits. An adaptive
  de-gain that protects against high-frequency oscillation / servo damage.
  \item \pp{PTCH2SRV\_RMAX} $=560/V$ --- manually scheduling the pitch \emph{rate
  limit} with speed to cap g-loading ($\pm 1$ g).
  \item D-term \& target filters (\pp{FLTD}, \pp{FLTT}) --- not speed-scheduled,
  but they keep the (speed-boosted) gains from amplifying noise at the low-speed
  end where scale is largest.
\end{itemize}

\subsection{What this means for your custom FDM}
\begin{enumerate}[leftmargin=1.6em]
  \item \textbf{One tune, whole envelope.} If your FDM computes moments as
  $\bar q\, S\, c\, C_m$ (i.e.\ moments scale with dynamic pressure), ArduPilot's
  $V^2$ scaling correctly covers 20--40 m/s from a single 30 m/s tune. This is a
  \textbf{validation test of the FDM}: tune at 30, fly at 20 and 40 --- if it
  stays stable, the aero model's speed dependence is physically consistent. If it
  oscillates at 40 but is fine at 30, either the FDM is not scaling moments with
  $\bar q$ correctly, or \pp{SCALING\_SPEED} is wrong.
  \item \textbf{uSTOL caveat --- blown wing / prop-wash.} Distributed propulsion
  augments the \emph{effective} dynamic pressure over the surfaces beyond
  freestream $V^2$. ArduPilot schedules only on airspeed, so in high-thrust /
  low-speed STOL regimes the real surface authority is higher than
  $(\text{\pp{SCALING\_SPEED}}/V)^2$ predicts --- the scheduler over-commands and
  you can get low-speed oscillation. Handles: tune (or AUTOTUNE at a low level,
  2--3) in the slow / high-blowing regime, or keep \pp{*\_RATE\_SMAX} set so the
  adaptive de-gain catches the over-command.
  \item \textbf{No per-throttle / per-flap table exists} in stock ArduPlane ---
  scheduling is airspeed-only. Thrust-dependent scheduling would be a model /
  firmware extension, not a parameter.
\end{enumerate}

% =====================================================================
\section{Part 4 --- TECS (height + speed), tune after pitch}
% =====================================================================

TECS coordinates \textbf{throttle and pitch} by managing total energy (potential
= height, kinetic = $\tfrac12 V^2$). Throttle adds/removes \emph{total} energy;
pitch \emph{shifts} energy between height and speed.

\textbf{Hard prerequisite:} the pitch rate loop must be solid first. About 90\% of
``TECS oscillation'' reports are actually a bad pitch loop one layer down.

\textbf{Setup order (measure in FBWA, then set):}
\begin{enumerate}[leftmargin=1.6em]
  \item \pp{THR\_MAX} --- enough to climb at \pp{PTCH\_LIM\_MAX\_DEG} at cruise
  (default 75\%).
  \item \pp{TRIM\_THROTTLE} --- throttle for level flight at cruise (from CTUN
  logs); also the airspeed-sensor-failure fallback.
  \item \pp{AIRSPEED\_MIN/MAX}, \pp{PTCH\_LIM\_MAX/MIN\_DEG}.
  \item \pp{TECS\_CLMB\_MAX} --- measured best climb rate at full throttle/cruise,
  set to 80\% of measured.
  \item \pp{TECS\_SINK\_MIN} --- sink at idle throttle, cruise speed.
  \item \pp{TECS\_SINK\_MAX} --- max safe descent without exceeding
  \pp{AIRSPEED\_MAX}.
\end{enumerate}

\textbf{Verify in a loiter:} holds altitude within approx 10 m, no porpoising.
Climb test: throttle should settle approx 80\% of MAX (not pinned), speed held,
pitch approx 5 deg below limit. If speed sags with throttle pinned: lower
\pp{TECS\_CLMB\_MAX} or raise \pp{THR\_MAX}.

\textbf{Fine-tune knobs (only if needed):}

\begin{longtable}{@{}p{3.4cm}p{2.0cm}p{9.4cm}@{}}
\toprule
\textbf{Param} & \textbf{Default} & \textbf{Effect}\\
\midrule
\endhead
\pp{TECS\_TIME\_CONST} & approx 5 s & Overall response. +1 to damp oscillation (max approx 10; if you need $>10$, the pitch loop is the real problem)\\
\pp{TECS\_PTCH\_DAMP} & 0.0 & +0.1 to damp height oscillation (do not exceed 0.5 without checking nav\_pitch noise)\\
\pp{TECS\_THR\_DAMP} & --- & Throttle-loop damping (needs airspeed sensor)\\
\pp{TECS\_INTEG\_GAIN} & 0.05--0.15 & How fast offsets trim out\\
\pp{TECS\_RLL2THR} & approx 10 & Extra throttle in turns (approx $10\times$ the m/s sink in a 45 deg bank). Low-aspect / draggy $\rightarrow$ higher. A uSTOL is likely draggy $\rightarrow$ expect higher\\
\pp{TECS\_SPDWEIGHT} & 1.0 & 0 = pitch holds height (ignores speed); 2.0 = pitch holds speed (gliders). Needs airspeed sensor\\
\pp{TECS\_PTCH\_FF\_K} / \pp{\_V0} & --- & Pitch feedforward vs speed demand (gliders $-0.04$, draggy $-0.08$)\\
\bottomrule
\end{longtable}

% =====================================================================
\section{Part 5 --- Navigation (L1), tune last}
% =====================================================================

L1 turns crosstrack/heading error into a \textbf{bank-angle demand}.

\begin{longtable}{@{}p{3.0cm}p{1.8cm}p{10.0cm}@{}}
\toprule
\textbf{Param} & \textbf{Value} & \textbf{Meaning}\\
\midrule
\endhead
\pp{NAVL1\_PERIOD} & 15 & Lower = sharper/tighter; too low = weaving/tail-wag. Higher = gentle. Tune $-5$ if turns lag, $+1$--$2$ if it weaves\\
\pp{NAVL1\_DAMPING} & 0.75 & Fine adjust in 0.05 steps after PERIOD; floor approx 0.6\\
\pp{WP\_RADIUS} & --- & approx distance flown in 2 s at cruise (approx 60 m at 30 m/s). Small = fly through WP then turn; large = cut the corner\\
\pp{WP\_LOITER\_RAD} & --- & Loiter/RTL/GUIDED circle radius\\
\pp{ROLL\_LIMIT\_DEG} & 30 & Max bank L1 may command (\pp{LEVEL\_ROLL\_LIMIT}=5)\\
\bottomrule
\end{longtable}

\textbf{Method:} fly a small rectangular AUTO loop you can watch; adjust
\pp{NAVL1\_PERIOD} until corners are crisp without weaving. A value of 15 is on
the tight/aggressive side (default 17--20) --- appropriate for a responsive
uSTOL, but watch for tail-wag.

% =====================================================================
\section{Part 6 --- Cruise, airspeed calibration, ground steering}
% =====================================================================

\begin{itemize}[leftmargin=1.4em]
  \item \textbf{Cruise tuning:} \pp{TRIM\_THROTTLE} = avg cruise throttle (from
  CTUN); \pp{PTCH\_TRIM\_DEG} = the few degrees of nose-up the aircraft flies at
  (preferred over re-levelling the accel). \pp{AIRSPEED\_CRUISE} is your hands-off
  target.
  \item \textbf{Airspeed calibration:} \pp{ARSPD\_RATIO} (approx 2.0) converts
  pressure to speed; \pp{ARSPD\_AUTOCAL}=1 learns it over a 5-min flight \emph{with
  turns} (never on a long straight leg --- bad estimates can cause a crash). In
  SITL this is largely moot --- the FDM provides truth airspeed; just keep the
  envelope honest. On real hardware: cover the pitot at boot for the offset cal.
  \item \textbf{Ground steering} (matters for a STOL with real takeoff/landing
  rolls): \pp{STEER2SRV\_P} = the aircraft's turn-circle \emph{diameter in metres}
  (measure by pushing it in a full-rudder circle, approx 4 m typical),
  \pp{\_I}=0.1, \pp{\_D}=0.02, \pp{\_TCONST}=0.5 (planes), \pp{GROUND\_STEER\_ALT}
  =5 m (hand-off to aero control above this). Tune takeoff in FBWA until it tracks
  straight with no rudder input, then move to AUTO takeoff. For a fast uSTOL the
  \pp{STEER2SRV\_DRT*} derating params reduce steering authority as you accelerate
  down the runway.
\end{itemize}

% =====================================================================
\section{Part 7 --- Doing it in SITL with a moving FDM}
% =====================================================================

\begin{enumerate}[leftmargin=1.6em]
  \item \textbf{Treat every FDM change as ``retune required''.} Aero edits like
  \pp{Cmq} (pitch-rate damping) directly change what \pp{PTCH\_RATE\_D} should be;
  ``post-stall \pp{CL}'' changes the low-speed / high-alpha behaviour. Old PIDs go
  stale.
  \item \textbf{Loop:} edit FDM $\rightarrow$ rebuild SITL $\rightarrow$ AUTOTUNE
  (level 6) $\rightarrow$ save params $\rightarrow$ envelope check (20/30/40 m/s)
  $\rightarrow$ inspect logs.
  \item \textbf{Logs are your truth} (this is where the real learning is):
  \begin{itemize}[leftmargin=1.2em]
    \item \pp{ATRP} --- AUTOTUNE progress (Type 0=roll, 1=pitch), demanded vs
    achieved rate, P evolution.
    \item \pp{RATE} / \pp{PIDR},\pp{PIDP},\pp{PIDY} --- per-axis target vs achieved
    rate and the P/I/D/FF contributions. This is how you see whether FF carries the
    load (good) or P does (under-tuned FF).
    \item \pp{CTUN} --- throttle, nav\_pitch vs pitch, airspeed (TECS health).
    \item \pp{NTUN} --- L1 crosstrack, desired vs actual bank (nav health).
    \item \pp{ATT} --- demanded vs actual roll/pitch/yaw.
  \end{itemize}
  \item \textbf{Use SITL to do what you cannot risk on hardware:} sweep
  \pp{SCALING\_SPEED} deliberately and watch the tune break; push to stall to
  validate post-stall \pp{CL}; run Monte-Carlo perturbation sweeps
  (\pp{LAT\_SIM\_MonteCarlo.cpp}, \pp{monte\_carlo\_overrides.txt}) to check tune
  robustness across CG/mass dispersions.
\end{enumerate}

% =====================================================================
\section{Part 8 --- This aircraft's current state \& next steps}
% =====================================================================

The current \pp{Latest.parm} is already a coherent, AUTOTUNE-derived tune:

\begin{longtable}{@{}p{4.2cm}p{2.6cm}p{6.8cm}@{}}
\toprule
\textbf{Parameter} & \textbf{Value} & \textbf{Note}\\
\midrule
\endhead
\pp{AIRSPEED\_CRUISE} & 30 & cruise target\\
\pp{AIRSPEED\_MIN / MAX} & 20 / 40 & envelope; MAX approx $2\times$ cruise (good)\\
\pp{SCALING\_SPEED} & 30 & matches cruise (correct anchor)\\
\pp{ARSPD\_USE / TYPE} & 1 / 2 & airspeed enabled in SITL\\
\pp{RLL\_RATE\_P/I/D/FF} & 0.214/0.214/0.034/0.400 & I = P signature\\
\pp{PTCH\_RATE\_P/I/D/FF} & 0.516/0.660/0.020/0.660 & I = FF signature\\
\pp{RLL2SRV\_TCONST} & 0.5 & \\
\pp{PTCH2SRV\_TCONST} & 0.75 & \\
\pp{PTCH\_LIM\_MAX/MIN\_DEG} & +20 / $-20$ & \\
\pp{ROLL\_LIMIT\_DEG} & 30 & \\
\pp{NAVL1\_PERIOD / DAMPING} & 15 / 0.75 & tight nav\\
\pp{AUTOTUNE\_LEVEL / AXES} & 6 / 7 & all 3 axes at 75 deg/s\\
\bottomrule
\end{longtable}

\textbf{High-value next moves:}
\begin{enumerate}[leftmargin=1.6em]
  \item \textbf{Verify gain scheduling holds:} fly steady at 20, 30, 40 m/s and
  check \pp{RATE}/\pp{PIDR}/\pp{PIDP} for oscillation at the extremes. If 40
  oscillates, the FDM may not scale moments with $\bar q$ correctly, or the rate
  D term is too hot at high $\bar q$.
  \item \textbf{Re-AUTOTUNE after \pp{Cmq} / post-stall changes} (the pitch D term
  especially).
  \item \textbf{Confirm \pp{AIRSPEED\_MIN}=20 is $\geq 1.25\times$ the FDM's actual
  stall} (you now have a real stall to measure).
  \item \textbf{Set \pp{TECS\_RLL2THR} deliberately} (a draggy uSTOL likely wants
  $>10$) and check loiter energy in turns.
\end{enumerate}

% =====================================================================
\section{Part 9 --- Doing it WITH Mission Planner}
% =====================================================================

Mission Planner has four places that matter for tuning, plus two live tools.

\subsection{The CONFIG tab --- where parameters live}
Top menu \textbf{CONFIG} (older builds: CONFIG/TUNING). Left column = sub-screens:

\begin{longtable}{@{}p{3.4cm}p{6.6cm}p{4.4cm}@{}}
\toprule
\textbf{Sub-screen} & \textbf{What it gives you} & \textbf{Use it for}\\
\midrule
\endhead
Basic Tuning & Curated panel of attitude/nav params & Quick edits to headline gains\\
Standard Params & Commonly-adjusted params, with descriptions/dropdowns & TECS, airspeed, L1, cruise\\
Advanced Params & Everything advanced users touch & Filters, SMAX, scaling, edge params\\
Full Parameter List & Every parameter, flat, searchable (the workhorse) & Anything not on the panels\\
Full Parameter Tree & Same data, hierarchical by prefix & Browsing a family (all \pp{TECS\_})\\
\bottomrule
\end{longtable}

\textbf{Full Parameter List buttons:}
\begin{itemize}[leftmargin=1.4em]
  \item \textbf{Search} --- type \pp{RLL\_RATE}, \pp{PTCH\_RATE}, \pp{TECS},
  \pp{NAVL1}, \pp{SCALING}, \pp{ARSPD}, \pp{STEER2SRV} to jump to a family.
  \item \textbf{Write Params} --- pushes edits to the autopilot (edits are not live
  until written; changed cells show green).
  \item \textbf{Refresh Params} --- re-read from vehicle (after enabling a feature).
  \item \textbf{Save to File / Load from File} --- backup/restore (critical for the
  FDM workflow).
  \item \textbf{Compare Params} --- diff current-vs-file, cherry-pick which to
  apply. This is how you A/B two tunes.
  \item \textbf{Reset to Default} --- factory reset (forces recalibration).
\end{itemize}

\textbf{Where each thing lives:} rate PIDs, attitude, yaw, ground steering,
airspeed cal $\rightarrow$ Full Parameter List (search the family). Gain
scheduling (\pp{SCALING\_SPEED}, \pp{AIRSPEED\_*}), TECS, nav $\rightarrow$
Standard Params or Full List.

\subsection{AUTOTUNE through Mission Planner}
There is no Autotune button for Plane --- it is a flight mode:
\begin{enumerate}[leftmargin=1.6em]
  \item Flight Data $\rightarrow$ Actions tab (or flight-mode dropdown)
  $\rightarrow$ set mode \textbf{AUTOTUNE}, or assign a switch via CONFIG
  $\rightarrow$ Flight Modes. Set \pp{AUTOTUNE\_LEVEL} and \pp{AUTOTUNE\_AXES}
  first.
  \item Fly the doublets (Part 2.3). Watch the Messages tab / HUD for autotune
  progress and ``saved'' notifications.
  \item Land, disarm $\rightarrow$ gains auto-save. Immediately Full Parameter List
  $\rightarrow$ Save to File to snapshot.
  \item A/B old vs new in the air: assign \pp{RCx\_OPTION}=180 and flip the switch.
\end{enumerate}

\subsection{The real-time Tuning graph (your manual-tuning feedback loop)}
\begin{enumerate}[leftmargin=1.6em]
  \item Flight Data screen $\rightarrow$ bottom edge of the map, tick
  \textbf{Tuning} $\rightarrow$ a black scrolling graph opens above the map.
  \item Double-click in that black window $\rightarrow$ a checklist of fields
  $\rightarrow$ tick what to plot.
  \item Classic plots from the manual-tuning docs:
  \begin{itemize}[leftmargin=1.2em]
    \item \pp{roll} vs \pp{nav\_roll} --- demanded vs achieved bank; spot overshoot
    while raising \pp{RLL\_RATE\_P/D}.
    \item \pp{pitch} vs \pp{nav\_pitch} --- pitch loop health (and TECS pitch
    demand).
    \item \pp{roll\_speed} / \pp{pitch\_speed} (rad/s) --- verify \pp{RMAX} limits.
    \item lateral accel --- tune \pp{YAW2SRV\_SLIP}/\pp{DAMP} (want approx 0).
  \end{itemize}
  \item Status tab: double-click any telemetry value to add it to the graph.
\end{enumerate}

\subsection{In-flight ``tuning knob'' (live gain tuning)}
ArduPlane can adjust a chosen gain live with a spare transmitter knob, so you dial
P/D in flight and watch the Tuning graph respond.
\begin{itemize}[leftmargin=1.4em]
  \item The mechanism is the \pp{TUNE} parameter family: \pp{TUNE} selects
  \emph{what} the knob adjusts, a tuning channel carries the knob, and min/max set
  the range. The knob updates approx $3\times$/s.
  \item In MP: Full Parameter List $\rightarrow$ search \pp{TUNE} $\rightarrow$ set
  the selector, channel, and min/max. (Confirm exact names in your firmware via
  that search; they vary slightly by version.)
  \item Workflow: pick the gain, set a min/max bracketing the current value, fly,
  twist until it just oscillates, back off, land, read the value, Write it.
  \item Note: the ``Extended Tuning'' screen with the CH6 \pp{TUNE} dropdown is the
  \emph{Copter} screen. For Plane you use the Full Parameter List \pp{TUNE} params.
\end{itemize}

\subsection{Dataflash log review (the real evidence)}
\begin{itemize}[leftmargin=1.4em]
  \item Flight Data $\rightarrow$ DataFlash Logs $\rightarrow$ ``Download DataFlash
  Log Via Mavlink'' (or pull the \pp{.bin} from the SITL log dir) $\rightarrow$
  opens the log graph. The web UAV LogViewer does the same and can re-fly in 3D.
  \item Graph \pp{RATE}/\pp{PIDR}/\pp{PIDP}/\pp{PIDY}, \pp{ATT}, \pp{CTUN},
  \pp{NTUN}, \pp{ATRP} (same messages as Part 7).
  \item Gain-scheduling check: fly steady at 20/30/40 m/s, overlay \pp{RATE}
  target-vs-achieved at each speed. Clean at all three = scheduling holds and the
  FDM's moments scale with $\bar q$ physically.
\end{itemize}

\subsection{Connecting MP to your SITL}
SITL exposes MAVLink on TCP \pp{127.0.0.1:5760} (and UDP 14550). Pick TCP in MP's
connection dropdown $\rightarrow$ \pp{127.0.0.1:5760}. \textbf{WSL2 caveat:} if MP
runs on Windows and SITL in WSL2, \pp{127.0.0.1} will not reach it --- use the
WSL2 VM IP (\pp{ip addr} in WSL), or run SITL with
\pp{--out=udp:<windows-ip>:14550} and have MP listen on UDP. If you drive SITL from
MAVProxy you can \pp{output add} a port for MP and keep MAVProxy too.

% =====================================================================
\section{Part 10 --- Doing it WITHOUT Mission Planner (headless / MAVProxy)}
% =====================================================================

This is often the better fit for a SITL / custom-FDM loop because it scripts.

\subsection{Launch SITL}
\begin{lstlisting}
# from the ardupilot root; builds + runs SITL with the console + map
Tools/autotest/sim_vehicle.py -v ArduPlane --console --map

# load your tuned parameters at startup
Tools/autotest/sim_vehicle.py -v ArduPlane --add-param-file=Latest.parm
\end{lstlisting}
This drops you into the \textbf{MAVProxy} console (the \pp{STABILIZE>} prompt).

\subsection{Read / set / save parameters (replaces the Full Parameter List)}
\begin{lstlisting}
param show RLL_RATE_*          # show a family (wildcards allowed)
param set PTCH_RATE_D 0.018    # set one value (live, no separate "write" step)
param fetch                    # re-read all from the vehicle
param save mytune.parm         # save current params to a file
param load mytune.parm         # load params from a file
param diff other.parm          # compare current vs a file  (A/B two tunes)
\end{lstlisting}

\subsection{AUTOTUNE without Mission Planner}
\begin{lstlisting}
param set AUTOTUNE_LEVEL 6
param set AUTOTUNE_AXES 7
mode FBWA                      # take off / get airborne first
mode AUTOTUNE                  # then switch in
# fly doublets via your RC / joystick, or with rc overrides:
rc 1 1800 ; rc 1 1200 ; rc 1 1500     # roll doublet (channel 1), repeat 20x+
rc 2 1800 ; rc 2 1200 ; rc 2 1500     # pitch doublet (channel 2)
mode FBWA                      # exit AUTOTUNE -> FF and I are saved
param save tune_post.parm
\end{lstlisting}
Autotune progress / ``saved'' messages print straight into the MAVProxy console.

\subsection{Real-time graphs (replaces the MP Tuning window)}
\begin{lstlisting}
module load graph
graph ATT.Roll ATT.DesRoll              # achieved vs demanded roll
graph ATT.Pitch ATT.DesPitch
graph RATE.R RATE.RDes                  # rate loop: achieved vs target
graph CTUN.Airspeed CTUN.ThrOut         # TECS / cruise
\end{lstlisting}

\subsection{Live ``knob'' tuning by script}
Without a physical knob you simply step the gain from the console and watch the
graph --- the scripted equivalent of the transmitter knob:
\begin{lstlisting}
param set RLL_RATE_D 0.030
# observe graph; if it oscillates, back off:
param set RLL_RATE_D 0.020
\end{lstlisting}

\subsection{Log review (replaces MP DataFlash review)}
SITL writes \pp{.BIN} logs under \pp{logs/} (and \pp{ArduPlane/logs/}).
\begin{lstlisting}
MAVExplorer.py logs/00000001.BIN        # interactive graphing + FFT + map
mavgraph.py    logs/00000001.BIN RATE.R RATE.RDes
mavlogdump.py  logs/00000001.BIN --type CTUN     # dump a message type to text
\end{lstlisting}
In \pp{MAVExplorer} use \pp{graph RATE.P*}, \pp{graph CTUN.*}, \pp{graph NTUN.*},
\pp{graph ATRP.*} --- the same messages as Part 7.

\subsection{Scripted / automated tuning (ties into your Monte-Carlo tooling)}
Because everything above is text, you can drive it from \textbf{pymavlink} to run
gain sweeps automatically and score them against log metrics --- the natural
companion to \pp{LAT\_SIM\_MonteCarlo.cpp}. Sketch:
\begin{lstlisting}
from pymavlink import mavutil
m = mavutil.mavlink_connection('udp:127.0.0.1:14550')
m.wait_heartbeat()
# set a candidate gain
m.mav.param_set_send(m.target_system, m.target_component,
                     b'PTCH_RATE_D', 0.018, mavutil.mavlink.MAV_PARAM_TYPE_REAL32)
# ... fly a scripted mission, pull the log, compute overshoot / RMS error,
#     loop over candidate gains, keep the best.
\end{lstlisting}

\subsection{Quick equivalence table}
\begin{longtable}{@{}p{6.6cm}p{8.0cm}@{}}
\toprule
\textbf{Mission Planner} & \textbf{Headless equivalent}\\
\midrule
\endhead
Full Parameter List + Write & \pp{param set} / \pp{param show} (live)\\
Save to File / Load from File & \pp{param save} / \pp{param load}\\
Compare Params & \pp{param diff}\\
AUTOTUNE via Actions tab & \pp{mode AUTOTUNE} + \pp{rc} doublets\\
Real-time Tuning graph & MAVProxy \pp{graph} module\\
Transmitter \pp{TUNE} knob & scripted \pp{param set} steps\\
DataFlash log review & \pp{MAVExplorer.py} / \pp{mavgraph.py}\\
\bottomrule
\end{longtable}

% =====================================================================
\section{Sources (official ArduPilot documentation)}
% =====================================================================
\begin{itemize}[leftmargin=1.4em]
  \item \url{https://ardupilot.org/plane/docs/common-tuning.html}
  \item \url{https://ardupilot.org/plane/docs/tuning-quickstart.html}
  \item \url{https://ardupilot.org/plane/docs/automatic-tuning-with-autotune.html}
  \item \url{https://ardupilot.org/plane/docs/new-roll-and-pitch-tuning.html}
  \item \url{https://ardupilot.org/plane/docs/roll-pitch-controller-tuning.html}
  \item \url{https://ardupilot.org/plane/docs/navigation-tuning.html}
  \item \url{https://ardupilot.org/plane/docs/tuning-cruise.html}
  \item \url{https://ardupilot.org/plane/docs/tecs-total-energy-control-system-for-speed-height-tuning-guide.html}
  \item \url{https://ardupilot.org/plane/docs/calibrating-an-airspeed-sensor.html}
  \item \url{https://ardupilot.org/plane/docs/tuning-ground-steering-for-a-plane.html}
  \item \url{https://ardupilot.org/plane/docs/airspeed-parameters-setup.html}
  \item \url{https://ardupilot.org/planner/docs/mission-planner-configuration-and-tuning.html}
  \item \url{https://ardupilot.org/planner/docs/mission-planner-flight-data.html}
  \item \url{https://ardupilot.org/planner/docs/common-uavlogviewer.html}
\end{itemize}

\end{document}
